\chapter{Process Execution}

%TODO: remove the instructions

In this final phase, software developers take over to implement the proposed solution.

\begin{enumerate}
    \item Create an executable BPMN model for \textbf{Camunda~8} (reuse the one created in~\cref{sec:to-be}). Model it in Camunda Modeler (Desktop or Web Modeler) and run it on a Camunda~8 engine, either locally (Camunda~8 Run) or on a Camunda SaaS cluster. Follow the lab session on BPMN execution in Camunda, and the \href{https://docs.camunda.io/}{Camunda~8 documentation}.
    \begin{enumerate}
        \item Simplify the model where needed, but keep the main flow. Use \textbf{user tasks only}: no scripts, connectors or job workers. A step that a real system would automate becomes a user task.
        \item The application must support different user roles. Define these roles and assign BPMN activities accordingly, based on responsibilities within the domain (in Camunda~8, as the candidate groups or assignees of the user tasks).
        \item Each process step should include a form (a Camunda Form linked to the user task) with appropriate validations as per the specification. Be creative! While advanced UI is welcome, simple text fields are also acceptable.
        \item Record a short walkthrough video demonstrating the process with multiple simulated participants: each role completes its tasks in Tasklist, and the running instance is shown in Operate. (Maximum duration: 5 minutes)
    \end{enumerate}
    \item Present the solution to the customer. Prepare a short, compelling presentation showcasing the To-Be process supported by your software solution. Imagine you are pitching to a customer who has invested 100,000 EUR in the project. Use your own voice. (Maximum duration: 2 minutes)
\end{enumerate}

\section{Walkthrough Presentation}
%insert a link to the video. 

\section{Sales Presentation}
%insert a link to the video. 
